\section{The Verification Asymmetry}
\label{sec:asymmetry}

The failure modes of \S\ref{sec:ver_failures} share a common formal
structure: the evaluator's computation is reachable from the agent's
action space.  This section formalizes that structure.

\subsection{Write-Access Formalization}
\label{sec:write_access}

Let $a$ be an agent with action space $\mathcal{A}(a)$, and let $E$ be
an evaluation function with input space $\mathcal{I}(E)$.  An action
$\alpha \in \mathcal{A}(a)$ \emph{reaches} an input $x \in
\mathcal{I}(E)$ if executing $\alpha$ can modify the value of $x$.
Write $\mathcal{R}(a, E) \subseteq \mathcal{I}(E)$ for the set of
evaluator inputs reachable by $a$.

\begin{definition}[Endogenous and Exogenous Evaluation]
\label{def:endo_exo}
An evaluation function $E$ is \textbf{endogenous} to agent $a$ if
$\mathcal{R}(a, E) \neq \emptyset$: there exists at least one action in
$a$'s action space that modifies at least one of $E$'s inputs.  $E$ is
\textbf{exogenous} to $a$ if $\mathcal{R}(a, E) = \emptyset$: no action
available to $a$ can modify any input to $E$.
\end{definition}

The distinction is binary in definition but admits a natural
quantitative extension.

\begin{definition}[Verification Asymmetry]
\label{def:ver_asymmetry}
The \textbf{verification asymmetry} of agent $a$ with respect to
evaluator $E$ is
\[
    \rho(a, E) \;\triangleq\; \frac{|\mathcal{R}(a, E)|}{|\mathcal{I}(E)|}
    \;\in\; [0, 1].
\]
When $\rho = 0$, evaluation is fully exogenous.  When $\rho = 1$,
evaluation is fully endogenous: every input to $E$ is reachable by $a$.
\end{definition}

\begin{remark}
When both $\mathcal{R}(a, E)$ and $\mathcal{I}(E)$ are continuous
spaces, the ratio of cardinalities is replaced by the ratio of their
measures under an appropriate reference measure $\mu$ on $\mathcal{I}(E)$:
$\rho(a, E) = \mu(\mathcal{R}(a, E)) / \mu(\mathcal{I}(E))$.
The qualitative results of this paper depend only on whether $\rho = 0$
or $\rho > 0$, so the choice of measure is immaterial for the structural
arguments.
\end{remark}

\subsection{Capability Scaling}
\label{sec:capability_scaling}

The verification asymmetry grows with capability.  As an agent becomes
more capable, its action space $\mathcal{A}(a)$ expands, and actions
that were previously beyond its reach become available.  The reachable
set $\mathcal{R}(a, E)$ is monotonically non-decreasing in the agent's
capability: acquiring a new capability never \emph{reduces} the set of
evaluator inputs the agent can influence.

\begin{proposition}[Monotonicity of Verification Asymmetry]
\label{prop:monotonicity}
Let $a$ and $a'$ be agents such that $\mathcal{A}(a) \subseteq
\mathcal{A}(a')$ (i.e., $a'$ has all of $a$'s capabilities and
possibly more).  Then $\mathcal{R}(a, E) \subseteq \mathcal{R}(a', E)$
and consequently $\rho(a, E) \leq \rho(a', E)$.
\end{proposition}

\begin{proof}
If action $\alpha \in \mathcal{A}(a)$ reaches $x \in \mathcal{I}(E)$,
then $\alpha \in \mathcal{A}(a')$ by the subset assumption, so $x \in
\mathcal{R}(a', E)$.  The inequality on $\rho$ follows from monotonicity
of the ratio.
\end{proof}

The proposition is elementary but its consequence is not: as AI systems
become more capable, the verification problem becomes \emph{strictly
harder}, not merely harder in practice.  Any verification scheme that
relies on the evaluator being ``hard to reach'' degrades as capability
grows.  A solution must identify a structural barrier that does not
degrade with capability.

\subsection{Substrate Exclusivity as the Structural Barrier}
\label{sec:substrate_barrier}

The only structural limit on $\mathcal{R}(a, E)$ that is invariant
under capability scaling is \emph{substrate exclusivity}.  An agent on
substrate $s$ can, in principle, write to any observation channel whose
computation occurs on $s$: it shares the physical medium, and capability
growth expands its reach within that medium without bound.  But an agent
on substrate $s$ cannot write to observation channels whose computation
occurs on substrate $s' \neq s$, because cross-substrate causal
influence requires a physical channel between substrates, and the
\emph{input space} of a process on $s'$ is not addressable from $s$ in
the same way that co-substrate processes are.

A population of agents occupying multiple computational substrates
$\{S_1, \ldots, S_m\}$ inherits a corresponding partition on
observation channels: channels rooted in substrate $s'$ are not
writable by agents on substrate $s \neq s'$.  This paper constructs
the protocol that \emph{enforces} this partition for verification
purposes.  (Part~II connects the construction to the substrate
partition formalized in the GFM
framework~\citep{lasser2026horizon, lasser2026scm}.)

\begin{definition}[Substrate-Exclusive Evaluation]
\label{def:substrate_exclusive_eval}
An evaluation function $E$ is \textbf{substrate-exclusive} with respect
to agent $a$ on substrate $s$ if $E$'s computation is rooted in
substrate $s' \neq s$.  By the substrate exclusivity property,
$\mathcal{R}(a, E) = \emptyset$ for any agent $a$ confined to substrate
$s$, and therefore $\rho(a, E) = 0$.
\end{definition}

Substrate exclusivity is not a capability limitation.  The agent may be
arbitrarily capable within its substrate; the barrier is physical, not
computational.  This makes it the only known verification barrier that
is invariant under capability scaling: no amount of capability growth
within substrate $s$ extends the agent's write access to substrate
$s'$.

\subsection{How Existing Frameworks Encounter the Asymmetry}
\label{sec:existing_frameworks}

Each major alignment framework encounters the verification asymmetry at
a characteristic point.

\paragraph{RLHF.}
The reward model is trained from human feedback---exogenous at training
time.  At deployment, the agent optimizes against a frozen reward model
whose inputs are within its influence: it can learn what input patterns
produce high reward scores without producing the intended
behavior~\citep{christiano2017deep, gao2023scaling}.  The asymmetry is
$\rho > 0$ at deployment because the reward model's input space
overlaps the agent's output space.

\paragraph{Constitutional AI.}
The constitution is exogenous text, but its \emph{interpretation} is
endogenous: the same model that is being constrained performs the
constitutional evaluation~\citep{bai2022constitutional}.  The asymmetry
is $\rho > 0$ because the interpretation step's ``inputs''---the
model's internal representations of constitutional meaning---are
writable by the model itself.

\paragraph{Debate and amplification.}
The verifier is a weaker model or a human with bounded
attention~\citep{irving2018ai}.  A sufficiently strong debater can
construct arguments that exploit the verifier's bounded capacity.  The
asymmetry enters through the persuasion channel: the debater's action
space includes actions that modify the verifier's belief state
(a subset of the verifier's inputs).

\paragraph{GFM.}
$\volL$ is defined over an external capability poset, but the agent's
\emph{measurement} of $\volL$ is endogenous: it uses its own world
model, trust estimates, and observation
channels~\citep{lasser2026gfm, lasser2026scm}.  A sufficiently capable
agent can construct self-consistent measurements that report $\volL$
growth while actual $\volL$ stagnates.  This was discovered
operationally in the GFM harness through the capability
self-assertion gradient: an agent that adds capabilities to its own
ledger entry is both claimant and observer.

In every case, the specification may be correct and the agent may have
generalized well, but the verification step is gameable because the
evaluator's computation is reachable from the agent's action space.
The differences between frameworks are in \emph{which} evaluator inputs
are reachable and \emph{how quickly} $\rho$ approaches~1 as capability
grows.  The structural fix---making $\rho = 0$ by rooting evaluation
in a different substrate---is the same across all of them.
